\documentclass[10pt,a4paper]{amsart}
\usepackage[margin=19mm]{geometry}
\usepackage{amsmath,amssymb,mathtools}
\usepackage[scr=rsfs,cal=euler]{mathalfa}
\usepackage{xfrac}
\usepackage[unicode,hidelinks,bookmarks=false]{hyperref}
\newcommand{\ie}{i.e.\ }
\newcommand{\cf}{cf.\ }
\newcommand{\resp}{respectively\ }
\newcommand{\Subsection}{\subsection}
\providecommand{\nobreakdash}{\nobreak\hbox{-}}
\setlength{\emergencystretch}{2em}
\multlinegap0pt
\usepackage{bm}
\usepackage{bbm}
\usepackage{dsfont}
\usepackage{xspace}
\usepackage{cancel}
\usepackage{mathtools}
\usepackage{amsthm}
\makeatletter
\@ifundefined{theorem}{\newtheorem{theorem}{Theorem}}{}
\@ifundefined{definition}{\newtheorem{definition}[theorem]{Definition}}{}
\@ifundefined{proposition}{\newtheorem{proposition}[theorem]{Proposition}}{}
\@ifundefined{remark}{\newtheorem{remark}[theorem]{Remark}}{}
\makeatother
\title[Regularity thresholds for anomalous dissipation and related ]{Regularity thresholds for anomalous dissipation and related phenomena in passive scalars}
\author{Marco Bagnara, Daniel W. Boutros, Camillo De Lellis, Svitlana Mayboroda}
\date{Source: arXiv:2603.11466v3 (CC BY 4.0)}
\begin{document}
\maketitle

\noindent\emph{Source and licence.} Regularity thresholds for anomalous dissipation and related phenomena in passive scalars. Version arXiv:2603.11466v3 (CC BY 4.0), distributed under the Creative Commons Attribution 4.0 International licence, as recorded in the arXiv licence metadata for this exact version and shown on the abstract page https://arxiv.org/abs/2603.11466v3. Full rights notice: licenses/arxiv-2603.11466-CC-BY-4.0.txt.\par\smallskip

\noindent\textit{Editorial scope.} Counted unit: the Remark whose source label is \texttt{None} in the article (source0.tex lines 540--542, source label \texttt{None}), delivered together with the article's own complete original proof of it (source0.tex lines 544--624, one \texttt{proof} environment of 81 lines). No mathematics is written, repaired or re-derived: statement and proof are the article's own text line for line. Required context retained uncounted: e:parabolic (lines 107--114); e:transport (lines 194--201); p:no-dissipation (lines 464--470); d:Richardson-dispersion (lines 523--528); p:no-Richardson (lines 532--538). Every definition, display and label of those lines is retained unchanged. Omitted from this package and not redistributed: 1--106, 115--193, 202--463, 471--522, 529--531, 539--539, 543--543, 625--750. Nothing inside a retained excerpt is omitted or altered: every retained body is delivered exactly as the article's lines are written, so no omission receipt of any kind accompanies this package. Citations: the retained excerpts carry no citation command at all, so no bibliography entry of the article is transcribed and no reference is redistributed. Reference adaptation: none was needed and none was made; every reference inside the delivered unit resolves inside it, so no unresolved reference or citation is produced. Printed slips kept literal: no printed word, symbol, hypothesis or number of the counted statement or of its proof was altered. Layout and class: the article's own class is \texttt{amsart} with packages including tikz, a4wide, float, amsmath, xfrac, stmaryrd, mathrsfs, bm, amsthm, mathtools, yfonts, amssymb, color, esint; the wrapper is the lane's pinned \texttt{amsart} editorial wrapper at 10pt on A4, which restates the theorem-like environments the retained text needs and adds the article's own macro definitions that the retained text needs (none), with extra wrapper packages (bm, bbm, dsfont, xspace, cancel, mathtools). The delivered numbering is the unit's own. Assets: no retained span contains a figure environment, a graphics-inclusion command, a PDF-page inclusion, a table or a TikZ picture; no image file is copied into this package, the package declares no asset dependency and no image file is redistributed with it. Licence: version arXiv:2603.11466v3 of this article is distributed under the Creative Commons Attribution 4.0 International licence, as recorded in the arXiv licence metadata for this exact version and shown on the abstract page https://arxiv.org/abs/2603.11466v3. A full rights notice is delivered at licenses/arxiv-2603.11466-CC-BY-4.0.txt. Characters: every retained excerpt is pure ASCII, so the delivered engine, XeLaTeX with its default Latin Modern text font, reports no missing character.
\begin{align}\label{e:parabolic}
\left\{
\begin{array}{ll}
\partial_t \theta_\varepsilon + (v\cdot \nabla) \theta_\varepsilon = \varepsilon \Delta \theta_\varepsilon\\ \\
\theta_\varepsilon (\cdot, 0) = \theta_{in}
\end{array}
\right.\, ,
\end{align}

\begin{equation}\label{e:transport}
\left\{
\begin{array}{ll}
\partial_t \theta + (v\cdot \nabla )\theta = 0 \\ \\
\theta (\cdot, 0) = \theta_{in}
\end{array}
\right.
\end{equation}

\begin{theorem}\label{p:no-dissipation}
Consider a bounded divergence-free vector field $v: \mathbb T^d \times [0,T] \to \mathbb R^d$ which satisfies the renormalization property . Let $\theta_{in}$ be a bounded initial datum and for any $\varepsilon>0$ we denote by $\theta_\varepsilon$ the unique bounded solution of equation \eqref{e:parabolic}. Then we have
\begin{equation}\label{e:no-dissipation}
\lim_{\varepsilon \downarrow 0}\, \varepsilon \int_0^T \int_{\mathbb T^d} |\nabla \theta_\varepsilon|^2 (x,t)\, dx\, dt = 0
\end{equation}
and $\theta_\varepsilon$ converges strongly in $C([0,T], L^2 (\mathbb{T}^d))$ to the unique bounded weak solution of \eqref{e:transport}.
\end{theorem}

\begin{definition}\label{d:Richardson-dispersion}
Denote by $d$ the geodesic distance on the torus. We say that $v$ supports  weak Richardson dispersion if there is a positive number $\delta > 0$ and a set $E$ of positive measure with the property that
\[
\liminf_{\varepsilon\downarrow 0} \mathbb E[ d (X_1^\varepsilon (x), X_1^\varepsilon (y))^2] > \delta \quad\mbox{for all $(x,y) \in E\times E$.}
\]
\end{definition}

\begin{proposition}\label{p:no-Richardson}
Assume $v$ is a (not necessarily autonomous) divergence-free vector field which satisfies the DiPerna-Lions property.  Let $X^\varepsilon_t (x)$ be as in Definition \ref{d:Richardson-dispersion} and let $X_t (x)$ be the unique regular Lagrangian flow of $v$. Then, for every positive $t$, 
\[
\lim_{\varepsilon\downarrow 0} \int_{\mathbb T^d} {\mathbb E} [ d (X^\varepsilon_t (x), X_t (x))^2]\, dx = 0\, ,
\]
where $d (y,z)$ denotes the (geodesic) distance between the points $y$ and $z$ on $\mathbb T^d$ (endowed with the flat metric). In particular $v$ cannot support Richardson dispersion.
\end{proposition}

\begin{remark}
Some authors adopt, as definition of Richardson dispersion, the statement that the variance of the one-point distribution $X^\varepsilon_t (x)$ remains strictly positive as $\varepsilon\downarrow 0$ for a nontrivial set of initial points $x$. Clearly Proposition \ref{p:no-Richardson} rules this out as well for any $v$ which satisfies the DiPerna-Lions renormalization property. 
\end{remark}

\begin{proof} Given that 
\[
{\mathbb E} [ d (X^\varepsilon_t (x), X_t^\varepsilon (y))^2]
\leq {\mathbb E} [ d (X^\varepsilon_t (x), X_t (x))^2] + {\mathbb E} [ d (X^\varepsilon_t (y), X_t (y))^2]
+ d (X_t (x), X_t (y))^2\, ,
\]
the main statement of the proposition and the usual Lusin property of measurable maps implies immediately the conclusion that $v$ cannot support Richardson dispersion.

Regard $\mathbb T^d$ as $\mathbb R^d/\mathbb Z^d$ and consider the vector field $v$ as a periodic vector field on $\mathbb R^d$. Both the solution of the stochastic ODE $X^\varepsilon_t$ and the regular Lagrangian flow $X_t$ in $\mathbb T^d$ can then be derived obviously from the corresponding objects in $\mathbb R^d$, which we denote by $Y_T^\varepsilon$ and $Y_T$. Hence, in this context it suffices to show that 
\begin{equation}\label{e:target}
\lim_{\varepsilon\downarrow 0} \int_{\Omega} {\mathbb E} [\min\{ |Y_T^\varepsilon (x) - Y_T (x)|^2, \bar C\}]
\, dx = 0
\end{equation}
for every bounded $\Omega$, every positive time $T$, and the geometric constant $\bar C = ({\rm diam}\, (\mathbb T^d))^2$. In fact if $\Omega$ contains $[0,1]^d$, then we have the inequality $d (X_T^\varepsilon (x), X_T (x))^2 \leq \min \{ |Y_T^\varepsilon (x) - Y_T (x)|^2, \bar C\}$ for all $x\in [0,1]^d \subset \Omega$ and the desired conclusion follows immediately.

\medskip

We next observe that the conclusions about the strong convergence of solutions to the parabolic equation \eqref{e:parabolic} to the unique solution of the transport equation, apply in the context of $\mathbb R^d$ as well if we assume that the initial data is in $L^2\cap L^\infty$. In fact the argument of Theorem \ref{p:no-dissipation} applies verbatim. 

\medskip

Consider now a fixed time $T$ and the parabolic equation solved {\em backward} in time on $(-\infty, T]$
\begin{equation}\label{e:backward-parabolic}
\left\{
\begin{array}{ll}
\partial_t \theta_\varepsilon + (v\cdot \nabla) \theta_\varepsilon = - \varepsilon \Delta \theta_\varepsilon\\ \\
\theta_\varepsilon (\cdot, T) = f
\end{array}
\right.
\end{equation}
together with the transport equation (solved backward in time as well)
\begin{equation}\label{e:backward-transport}
\left\{
\begin{array}{ll}
\partial_t \theta + (v\cdot \nabla) \theta = 0 \\ \\
\theta (\cdot, T) = f\, .
\end{array}
\right.
\end{equation}
If $f \in L^2\cap L^\infty$, the arguments of Theorem \ref{p:no-dissipation} imply that 
\begin{itemize}
\item[(i)] $\|\theta_\varepsilon (\cdot, 0) - \theta (\cdot, 0)\|_{L^2}$ converges to $0$ as $\varepsilon\downarrow 0$.
\end{itemize}
On the other hand we also know that 
\begin{itemize}
\item[(ii)] $\theta (x,0) = f (Y_T (x))$ for a.e. $x$ by the DiPerna-Lions theory;
\item[(iii)] $\theta_\varepsilon (x,0) = {\mathbb E} [ f (Y_T^\varepsilon (x))]$ by the Feynman-Kac formula. 
\end{itemize}
Since $\Omega$ is bounded and $Y_T$ is the regular Lagrangian flow of the bounded vector field $v$, if we choose $R$ large enough we have $Y_T (\Omega) \subset B_R$. Consider now $f := \mathbf{1}_{B_R}$ and the corresponding solutions of \eqref{e:backward-parabolic} and \eqref{e:backward-transport}. Then $\theta_\varepsilon (\cdot, 0)$ converges strongly in $L^2$ to $\theta (\cdot, 0)$, which is identically equal to $1$ on $\Omega$ because of (ii) and $Y_T (\Omega) \subset B_R$. But then $1-\theta_\varepsilon (\cdot, 0)$ converges strongly in $L^2 (\Omega)$ to $0$. Using the Feynman-Kac formula, this implies
\begin{equation}\label{e:escaping}
\lim_{\varepsilon \downarrow 0} \int_\Omega \mathbb E[1-\mathbf{1}_{B_R} (Y_T^\varepsilon (x))]\, dx = 0\, .
\end{equation}
Clearly we have 
\begin{align*}
|Y^\varepsilon_T (x) - Y_T (x)|^2
\leq & 2 |Y^\varepsilon_T (x) \mathbf{1}_{B_R} (Y^\varepsilon_T (x)) - Y_T (x)|^2
+ 2 |Y^\varepsilon_T (x)|^2 (1-\mathbf{1}_{B_R} (Y^\varepsilon_T (x)))\, , 
\end{align*}
which in turn implies 
\begin{align*}
\min \{|Y^\varepsilon_T (x) - Y_T (x)|^2, \bar C\}
\leq 2  |Y^\varepsilon_T (x) \mathbf{1}_{B_R} (Y^\varepsilon_T (x)) - Y_T (x)|^2
+ \bar C (1-\mathbf{1}_{B_R} (Y^\varepsilon_T (x))\, .
\end{align*}
In particular, from \eqref{e:escaping} we conclude that, in order to show \eqref{e:target}, it suffices to show that the integral
\begin{align*}
&\int_\Omega {\mathbb E} [|Y_T^\varepsilon (x) \mathbf{1}_{B_R} (Y_T^\varepsilon (x)) - Y_T (x) |^2]\, dx \\
= & \underbrace{\int_\Omega {\mathbb E} [|Y_T^\varepsilon (x) \mathbf{1}_{B_R} (Y_T^\varepsilon (x)) - Y_T (x)\mathbf{1}_{B_R} (Y_T (x)) |^2]\, dx}_{=: I (\varepsilon)}
\end{align*}
vanishes as $\varepsilon\downarrow 0$. 

Expanding the square and using linearity of the expectation we can write 
\begin{align*}
I (\varepsilon) &= \int_\Omega {\mathbb E} [|Y_T^\varepsilon (x)|^2 \mathbf{1}_{B_R} (Y_T^\varepsilon (x))]\, dx\\
& - 2 \int_\Omega Y_T (x) \mathbf{1}_{B_R} (Y_T (x)) \cdot {\mathbb E} [Y_T^\varepsilon (x) \mathbf{1}_{B_R} (Y_T^\varepsilon (x))]\, dx\\
&+ \int_\Omega |Y_T (x)|^2\mathbf{1}_{B_R} (Y_T (x)) \, dx\, . 
\end{align*}
Consider now the equation \eqref{e:backward-parabolic} with $f (x) = |x|^2 \mathbf{1}_{B_R} (x)$. Its solution is then given, at time $0$, exactly by the integrand of the first integral in the expression above. For $\varepsilon\downarrow 0$ we know that this converges, strongly in $L^2 (\mathbb R^d)$, to the unique solution, at time $0$, of \eqref{e:backward-transport} with the same initial data. This is then $|Y_T (x)|^2 \mathbf{1}_{B_R} (Y_T (x))$, namely the integrand of the third integral. But then, since $\Omega$ is bounded, the first integral converges, as $\varepsilon\downarrow 0$, to the third integral. 

On the other hand for the second integral we can consider the (vector-valued) initial data $f (x) = x \mathbf{1}_{B_R} (x)$ and again argue that ${\mathbb E}  [Y_T^\varepsilon (x) \mathbf{1}_{B_R} (Y_T^\varepsilon (x))]$ is the solution of the corresponding backward parabolic equation at time $0$. Again the latter converges strongly in $L^2 (\mathbb R^d)$ to the solution of the backward transport equation at time $0$, which is $Y_T (x) \mathbf{1}_{B_R} (x)$. So the second integral is also converging to the third integral. Given however the factor $-2$ in front of it, in the limit as $\varepsilon\downarrow 0$ the three expression cancel out, showing that $I (\varepsilon)\downarrow 0$. 
\end{proof}

\end{document}
